\chapter{Introduction}\label{ch:intro}

\section{Motivation}
Training a neural network is expensive, and a large body of work proposes ways to make it cheaper: better optimisers, curricula, sparsification, selection of examples, and alternatives to the backpropagation algorithm itself. Each new proposal usually arrives with a plot showing it beating a baseline. Careful evaluations repeatedly find that such plots overstate the gain. When optimisers are compared with a tuned baseline and a stated protocol, the ranking is less clear than the individual papers suggest \citep{schmidt2021crowded}. When several popular efficient-training algorithms for language models are compared with a baseline trained under a fully decayed learning-rate schedule and an equal budget, their gains largely disappear \citep{kaddour2023notrain}. Different cost measures (parameters, floating-point operations, wall-clock time) can also disagree about which method is cheapest \citep{dehghani2022efficiency}.

Three problems recur. First, the baseline is tuned less than the new method, or tuned on too few random seeds to be reliable. Second, the cost of the extra machinery (additional forward passes, a second network, selection steps) is not counted. Third, the question ``did it work?'' is settled after looking at the data, so that the threshold, the metric and the comparison group can drift towards the result.

This report is a small, controlled case study of how to avoid these problems, and of what happens to two concrete ideas when they are tested that way. It is deliberately small: the model is a four-layer multilayer perceptron on MNIST. The smallness is the point. A tiny network can be trained hundreds of times on a laptop CPU, so every arm can be tuned with the same budget, every comparison can use many paired seeds, and each hypothesis can be written down, locked in version control and then tested once.

\section{The two ideas}
\paragraph{Idea 1: a meeting loss.}
Target propagation is a family of credit-assignment methods that replace the backward gradient pass by learned \emph{inverse} networks that map a desired output back to a desired hidden state \citep{bengio2014auto,lee2015dtp}. They are attractive because they avoid the weight-transport problem of backpropagation, but they are known to scale poorly \citep{bartunov2018assessing,meulemans2020theoretical}. The first idea of this report, which gives the repository its name, uses a target-propagation signal not as a replacement for backpropagation but as an \emph{auxiliary} term next to it: at the middle layer $h$ of the network, an extra loss pulls $h$ towards a difference-corrected target computed by a learned inverse network, so that a signal from ``below'' and the usual gradient from ``above'' meet. The hope was that the extra signal would reduce the FLOPs needed to reach a given accuracy.

\paragraph{Idea 2: a selective backward pass.}
While analysing the gradients of the baseline network after Study~1, a simple measurement stood out: once training is under way, a small share of the examples carries nearly all of the per-example gradient norm (the top $30\%$ of examples by loss hold $98\%$ of it by step 500; \cref{sec:concentration}), and the meeting gradient turned out to live on the same rows. This suggests running the forward pass on the whole batch (which also yields every example's loss) and the backward pass only on the highest-loss rows. This is the idea behind online hard example mining \citep{shrivastava2016ohem}, online batch selection \citep{loshchilov2016online}, Ordered SGD \citep{kawaguchi2020ordered} and, in its closest form, Selective-Backprop \citep{jiang2019selective}. The report does not claim the idea is new. It asks a narrower question: in a setting where cost is counted exactly, against baselines tuned as carefully as the method, how much does it save, why, and which design choices matter?

\section{Research questions}
\begin{description}[leftmargin=2.6em,style=nextline]
\item[RQ1] Does the meeting loss reduce the FLOPs needed to reach a fixed validation accuracy, compared with an equally well tuned baseline and with a deep-supervision control that has the same cost?
\item[RQ2] If it does not, why not? Is the meeting signal wrong, too weak, or redundant with backpropagation?
\item[RQ3] Does a selective backward pass reduce FLOPs-to-target, and is the reduction more than an effect of using a smaller effective batch?
\item[RQ4] Which design choices of the selective pass matter, and can a lean variant beat a tuned reimplementation of Selective-Backprop when all extra forward passes are charged?
\item[RQ5] How reliable is the usual practice of tuning a baseline on two seeds and comparing on a handful of others, and how well does a FLOPs saving translate into time?
\end{description}

\section{Contributions}
The contributions are empirical and methodological. None of them is a new algorithm.
\begin{enumerate}
\item \textbf{A pre-registered negative result} for the meeting loss (\cref{ch:study1}): on 20 paired seeds the loss costs about $36\%$ more FLOPs than a tuned baseline, and a same-cost deep-supervision control does not help. Seven further variants and controls, plus three mechanism checks (\cref{ch:mechanism}), narrows down the explanation without settling it.
\item \textbf{A pre-registered positive result} for the selective backward pass (\cref{ch:study2}): about $35$--$38\%$ fewer FLOPs than a tuned baseline and $30$--$36\%$ fewer than a baseline tuned for a smaller batch.
\item \textbf{A pre-registered improvement} (\cref{ch:study3}): removing the warm-up and keeping fewer rows saves a further $19$--$22\%$ against the first variant and about half the FLOPs of the baseline, and it beats a tuned Selective-Backprop reimplementation by $47$--$54\%$ in FLOPs.
\item \textbf{A protocol and its lessons} (\cref{ch:protocol,ch:discussion}): a crossing rule for FLOPs-to-target, a noise-floor check by re-tuning the baseline, evidence tiers, and a record of every post-hoc decision (\cref{app:deviations}). The central lesson is that tuning on two seeds is unreliable: two tunings of the same baseline that tied at selection differed by up to $34\%$ on fresh seeds.
\item \textbf{A full audit trail}: every number in this report is recomputed from the raw logs of the repository by scripts shipped with the report source.
\end{enumerate}

\section{Scope and non-goals}
All experiments use one architecture (784--256--128--256--10), one dataset (MNIST), one optimiser (AdamW), batch size 128 and a horizon of 3\,000 steps on a CPU. Cost is counted in matrix-multiplication FLOPs, not measured as time. The report does \emph{not} claim that either idea works on deep convolutional networks, language models or GPUs, and it does not claim that its Selective-Backprop reimplementation reproduces the published wall-clock gains of that method. \Cref{sec:limitations} lists the limitations in full.

\section{Evidence tiers and how to read the results}
Because the work was done in several rounds, the report labels every result with its evidential status:
\begin{itemize}
\item \tierC{} results test a hypothesis, margin and decision rule that were written down and locked in version control \emph{before} the confirmatory runs (\cref{sec:prereg}). Only these results support claims.
\item \tierE{} results come from looks at the data that were not pre-registered: pilots, mechanism checks, headroom studies. They generate hypotheses. They are reported with intervals, but without verdicts.
\item \tierV{} results reproduce an earlier number to check that the code still behaves the same.
\end{itemize}

\section{Structure of the report}
\Cref{ch:background} reviews related work. \Cref{ch:methods} defines the model, the two ideas and their cost. \Cref{ch:protocol} describes the experimental protocol. \Cref{ch:study1,ch:mechanism} report Study~1 and the mechanism checks; \cref{ch:study2,ch:study3} report Studies~2 and~3. \Cref{ch:discussion} discusses what was learned, including threats to validity, and \cref{ch:conclusion} concludes with future work. The appendices contain the reproduction guide, hyper-parameters, the hypothesis register, exploratory tables and a list of protocol deviations. \Cref{fig:overview} shows the overall design.

\begin{figure}[t]
\centering
\resizebox{\textwidth}{!}{% Project timeline (Chapter 4). Only dated events: commit dates, tag dates and the dates in the run-log file names.
\begin{tikzpicture}[x=1.45cm,
  ev/.style={draw,rounded corners=2pt,font=\scriptsize,align=center,fill=white,inner sep=3pt,text width=2.3cm},
  lock/.style={ev,fill=oiBlue!12,draw=oiBlue},
  expl/.style={ev,fill=oiOrange!12,draw=oiOrange!80!black}]
\draw[-{Stealth},thick] (-0.3,0) -- (10.6,0);
\foreach \x/\d in {0/20,1/21,2/22,3/23,4/24,5/25,6/26,7/27,8/28,9/29,10/30}{
  \draw (\x,0.06)--(\x,-0.06); \node[font=\scriptsize,below] at (\x,-0.08) {\d};}
\node[font=\scriptsize,anchor=east] at (-0.4,0) {Oct 2026};
\node[expl] (a) at (0.9,1.45) {first commit (20th);\\pilots, diagnostics\\(21st--22nd)};
\node[lock] (b) at (3.0,2.15) {\textbf{prereg-v1} locked\\\code{3378d36}\\H1, H6: 150 runs};
\node[lock] (d) at (6.4,2.15) {\textbf{prereg-v2} locked\\\code{cc2830c}\\H7a, H7b: 220 runs (27th)};
\node[lock] (g) at (9.55,1.45) {\textbf{prereg-v3} locked\\\code{a49e9e3}\\H8--H10: 260 runs};
\node[expl] (e) at (6.0,-1.45) {dose-response\\runs (26th)};
\node[expl] (f1) at (8.1,-1.45) {headroom: 150 runs (28th)};
\node[expl] (f2) at (9.9,-1.45) {look 2: 120 runs (29th)};
\draw[black!50] (a.south)--(1,0); \draw[black!50] (b.south)--(3,0); \draw[black!50] (d.south)--(6,0); \draw[black!50] (g.south)--(10,0);
\draw[black!50] (e.north)--(6,0); \draw[black!50] (f1.north)--(8,0); \draw[black!50] (f2.north)--(9,0);
\end{tikzpicture}
}
\caption{The project in time. Blue boxes are pre-registration lock commits (git tags), orange boxes are exploratory work; dates are those of the commits and of the run-log files. The whole project, from first commit to the last confirmatory result, took eleven calendar days on a single laptop.}
\label{fig:overview}
\end{figure}
